\documentclass{beamer}

% Packages (if needed)
\usepackage{graphicx}

% Presentation information
\title{Your Presentation Title}
\subtitle{Optional Subtitle}
\author{Your Name}
\institute{Your University/Institution}
\date{\today}

\begin{document}
%%%%%%%%%%%%%%%%%%%%%%%%%%%%%%%%%%%%%%
%%%%%%%%%%%%%%%%%%%%%%%%%%%%%%%%%%%%%%


% Title slide
\frame{\titlepage}


%%%%%%%%%%%%%%%% Begin Standard solvers
\begin{frame}
\frametitle{Grid based solvers}
- FDM
- FEM
cg, 
- 

Challanges in high dimensions
- need to store
How to address them
- parallelize

Hardware requirements
- RAM memory
- SSD storage
\end{frame}
%%%%%%%%%%%%%%%% End Standard solvers


%%%%%%%%%%%%%%%% Begin PINNs
\begin{frame}
\frametitle{Physics Informed Neural Networks (PINNs)}
Intro
- show diagram
- loss from NN
- how to train?
- sample points

Challanges in high dimensions
- laplace operator computation

How to handle them
- stoch laplace comp
- domain - based on loss
- parallel


Advantages / Disagvantages
\begin{itemize}
    \item flexibility
    \item memory efficient
\end{itemize}

Hardware requirements
- GPU(s)
\end{frame}
%%%%%%%%%%%%%%%% End PINNs


%%%%%%%%%%%%%%%% Begin deep BSDE
\begin{frame}
\frametitle{Deep BSDE method}
- just mention as next to try
\end{frame}
%%%%%%%%%%%%%%%% End deep BSDE





%%%%%%%%%%%%%%%%%%%%%%%%%%%%%%%%%%%%%%
%%%%%%%%%%%%%%%%%%%%%%%%%%%%%%%%%%%%%%
\end{document}
